
 \chapter{Nombres relatifs}
  \begin{itemize}
  \it Peut-on mesurer à l'aide des nombres relatifs les 
  grandeurs suivantes : fortune d'une personne, gains d'un
  joueur, vitesse d'un mobile sur un axe, date d'un évènement ?
  \it Placer sur un axe $(xy)$ les points suivants donnés par leurs abscisses, exprimées en centimètres : $A(+3); B(-2); C(-5); D(+4)$.\\En déduire les valeurs de $\overline{AB}$, $\overline{AC}$, $\overline{AD}$, $\overline{BC}$, $\overline{BD}$ et $\overline{CD}$.
  \it Sur un axe $(x'x)$ d'origine $O$, on construit les points $A(+3)$, $B(+11)$, et $M(+7)$. Calculer les mesures algébriques $\overline{AB}$, $\overline{AM}$, $\overline{MB}$ et $\overline{OM}$. Que représente le point $M$ pour le segment $[AB]$ et que peut-on dire des vecteurs $\overrightarrow{MA}$ et $\overrightarrow{MB}$ ? 
  \it On construit sur un axe $(x'Ox)$ les points 
  $A(+5), B(+9), C(-1)$ et $D(-5)$. Comparer les vecteurs $\overrightarrow{AB}$ et $\overrightarrow{DC}$ ainsi que les vecteurs $\overrightarrow{AD}$ et $\overrightarrow{BC}$. Que représente le point $M(+2)$ pour chacun des segments $[AC]$ et $[BD]$ ? 
  \it On place sur un axe $(xy)$ deux points $A$ et $B$ d'abscisses $+1$ et $-3$. Quelles sont les abscisses des points qui partagent le segment $[AB]$ en huit parties égales ? 
  \it Soient $A$ et $B$ deux points d'abscisses $-2$ et $-5$. Quelle est l'abscisse du milieu $M$ de $[AB]$ ? Quelle est l'abscisse du point $N$ tel que $\overline{BN} = 2\overline{AN}$ ?
  \it Soit un point d'abscisse $+5$ ; quelle est l'abscisse du point $A'$ symétrique de $A$ par rapport à l'origine $O$ des abscisses ? Soit de même $B$ d'abscisse $-1$ et $B'$ son symétrique par rapport à $O$. Évaluer $\overline{AB}$ et $\overline{A'B'}$. 
  \it Dans l'échelle d'un thermomètre Fahrenheit la division $32$ correspond au $0^o$ centigrade et la division $212$ à $100^o$ centigrades. 
  \begin{enumerate}
  \item À combien de degrés centigrades correspond un degré Fahrenheit ? 
  \item Trouver les températures centigrades correspondant à $+50$ degrés Fahrenheit et à $+5$ degrés Fahrenheit. 
  \end{enumerate}
  \it \begin{enumerate}
  \item À combien de degrés Fahrenheit correspond un degré centigrade ?
  \item Trouver la température Fahrenheit correspondant à
  $+15^o$ centigrades et à $-15^o$ centigrades.
  \end{enumerate}
  \end{itemize}